\documentclass[11pt]{article}
\usepackage[T1]{fontenc}
\usepackage{lmodern}
\usepackage[margin=0.85in]{geometry}
\PassOptionsToPackage{obeyspaces,spaces}{url}
\usepackage{amsmath,amssymb,mathtools,booktabs,longtable,array,enumitem,microtype,xcolor,xurl,hyperref,fancyhdr}
\hypersetup{colorlinks=true,linkcolor=blue!45!black,urlcolor=blue!45!black,
 pdftitle={Initial draft assessment: Comparing the October and November Yang--Mills submissions},
 pdfauthor={GPT Astra under human supervision},
 pdfsubject={Comparison of the common existence results; perturbative matching deferred}}
\setlength{\parindent}{0pt}
\setlength{\parskip}{5pt plus 1pt}
\setlength{\headheight}{14pt}
\setlist{leftmargin=1.5em,itemsep=2pt,topsep=3pt}
\pagestyle{fancy}\fancyhf{}
\fancyhead[L]{\small October--November Yang--Mills comparison}
\fancyhead[R]{\small 25 November 2026}
\fancyfoot[C]{\thepage}
\newcommand{\R}{\mathbb R}
\newcommand{\C}{\mathbb C}
\newcommand{\N}{\mathbb N}
\newcommand{\g}{\mathfrak g}
\newcommand{\code}[1]{{\small\nolinkurl{#1}}}
\newcommand{\src}[1]{\textsf{[#1]}}
\newcommand{\oldlink}[2]{\href{https://github.com/yang-mills-group/OpenAI-YM-20261028/blob/3456e5e75cf8d4cd558955a5c7638fba2d68f464/#1}{#2}}
\newcommand{\reviewlink}[2]{\href{https://github.com/yang-mills-group/YM/blob/55ba9fc36293b3f4d52c6b1f813169a0be76ee98/submissions/openai-20261028/#1}{#2}}
\begin{document}
\begin{center}
{\large\bfseries Initial draft assessment}\par
{\large Comparing the October and November Yang--Mills submissions}\par
{\small Common construction and strengthened existence checks}\par
{\small AI generated by GPT Astra under human supervision}\par
{\small 25 November 2026}
\end{center}

\textbf{Scope.} This document compares the existence material shared by the two
submissions, with particular attention to the mathematical requirements in
OpenAI's November existence Comparator target. It also compares that target
with the group's stronger October acceptance theorem. The new perturbative
matching theorem, its diagrammatic reference, and its positive-order
correlation and OPE calculations are reserved for a later review.

\textbf{Current finding.} November retains a large common construction core
and makes substantially more of its relationship to Yang--Mills explicit in
the public existence statement. The most consequential additions connect
microscopic fields to reconstructed fields, identify phases through moments
of the same Wilson hierarchy, require a Ward-normalized stress tensor, and
impose a free-Maxwell short-distance comparison. Our October acceptance
theorem already retains much of the construction data missing from the small
October target. Nevertheless, equivalence between our October target and
November's target has not been proved. The latter's explicit Maxwell
comparison is stronger in its radius and phase quantifiers. Section~\ref{sec:proofs}
traces how these obligations are proved and distinguishes additional
arguments from new public presentations of existing results.

\textbf{Verification status.} The November Existence and Matching targets
both passed Comparator, and both subsequently passed a repeat with Comparator
\code{d03acab154d269c06e60e4de7e4cc85deebff94b} and the Lean 4.34.0 checking
kernel. The final repeat finished on 10 November at 00:19 Pacific time.
These are completed receipts, not estimates of a running process. The runs
used macOS development mode without Linux/Landrun isolation or an independent
second kernel. Passing Matching is recorded as a mechanical fact; its physical
interpretation is outside this document. Section~\ref{sec:verification}
gives the precise scope and records.

The group's October checks have now also been transferred to the November
construction. Both combined targets passed Comparator on 11 November, with
only the three standard Lean axioms. The new theorem retains one native
construction witness for November's public existence conditions and our
independent checks; it does not assert equivalence of the two submissions'
chosen theories. Section~\ref{sec:remaining} describes this additional result.
\src{G3,E4}

\section{Three statements that must be distinguished}
\label{sec:three}

The October public target, the group's October acceptance target, and the
November public target ask different mathematical questions. A change in
the public target need not mean that the older construction lacked a proof
of the newly exposed property. It can mean that the older short statement
discarded information established by its proof.

\emph{October's public target} is \code{YangMills.JointExistence}, proved by
\code{YangMills.jointExistence}. Its final hierarchy satisfies four named
properties: Wilson convergence, reconstruction, an OPE with estimates, and a
nonfree trace test. This short interface leaves the microscopic insertions
and some relations between the resulting structures comparatively free.
\src{O1}

\emph{The group's October target},
\code{YMReview.Acceptance.combined_acceptance}, retains an actual native
\code{Family} and a \code{YM.Spec.Final} witness. Its clauses use that same
witness for the independent lattice comparison, scalar OS axioms,
translation Ward identity, operator content, leading trace law, and positive
continuum probe. Its \code{SourceConclusions} also retains stronger submitted
reconstruction, OPE, contact and short-distance results. Retaining these
source predicates is useful, but does not turn every internal definition into
an independently validated physical specification. \src{G1}

\emph{November's public target} has the following form:
\[
 \forall g\;\forall h:g.\mathrm{IsSimple},\quad
 \exists B:\mathrm{Maxwell.Reference}(g),\quad
 \exists R:\mathrm{Wilson.Construction}(g),\quad R.\mathrm{Realizes}(B).
\]
Here \code{Construction} contains one Wilson family and a function selecting
a refinement and hierarchy for every requested increasing volume sequence.
The content of \code{Realizes}, rather than the short theorem displayed in
the challenge file, is the central object to review. Its proof has its own
import closure and does not import \code{Comparators.Matching.Solution}.
It does share substantial internal analysis with other parts of the release.
\src{N1,N2}

\begingroup\small\setlength{\tabcolsep}{4pt}
\renewcommand{\arraystretch}{1.15}
\begin{longtable}{@{}>{\raggedright\arraybackslash}p{.16\textwidth}>{\raggedright\arraybackslash}p{.25\textwidth}>{\raggedright\arraybackslash}p{.27\textwidth}>{\raggedright\arraybackslash}p{.25\textwidth}@{}}
\toprule
Question & October public target & Group's October target & November Existence target\\\midrule\endfirsthead
\toprule Question & October public target & Group's October target & November Existence target\\\midrule\endhead
Lattice fields & Continuous, local, gauge-invariant insertions. & Native polynomial sections, classical jet normalization and RG transfer certificates. & Linear graded raw fields, unit, conjugation and OS adjoint, covering the physical field space.\\
Bare coupling & A positive logarithmic slope. & Retains the native calibrated family; a separate lattice-law comparison fixes conventions. & Explicit one-loop slope and a common cutoff/coupling trajectory.\\
Limits & Full UV limit at fixed volume; an IR subsequence for every requested volume sequence. & Same pattern through the retained native witness. & Also exposes a common finite-volume point table and local uniform point convergence.\\
Phases & A mixture identity; reconstruction attached to the selected theory. & Native spectral phase certificate and all-phase results on the same witness. & Every phase reconstructed; separated word trains fix joint moments of phase readings.\\
Stress trace & A supplied grade-four field called \code{centeredTrace}. & Identified with the trace of the same tensor checked by the translation Ward test. & Defined from the supplied stress tensor; translation and rotation Ward laws required.\\
Maxwell limit & Absent from the short target. & Explicit \code{FirstExtension}: mean hierarchy, common shrinking sequence, moving frames, all derivatives on separated compacta. & Mean and every phase, all sufficiently small real radii, frame bounds and fixed-field identification.\\
OPE & Coefficient and finite-truncation error bounds. & Retains these and native nested results; a separate group Comparator checks OPE alignment. & Also exposes coefficient algebra and finite-tree power estimates.\\
Independent OS test & Not the group's formulation. & Compiled scalar bridge, including E0$'$ growth. & Group's bridge now proved for the November witness, including identification with its public selected trace.\\
IR dynamics & No unique infinite-volume limit or mass gap. & No such additional conclusion. & No such additional conclusion.\\
\bottomrule
\end{longtable}\endgroup

\section{The common construction: inputs and order of limits}
\label{sec:limits}

Both submissions start with a compact matrix gauge group whose real Lie
algebra is nonabelian and simple, and with the Wilson lattice measure.
The compared public definitions of the gauge group and simplicity condition
are unchanged after removal of comments and whitespace. The same is true
of the compared base Wightman definitions. The Wilson action is still the
specified plaquette action, not an unspecified measure selected after the
continuum functions have been found. \src{O1,N3,E2}

A member of the regulated family requires more than a lattice size. It
includes an inverse spacing $D_N$, a physical period $L_p$, the Wilson
parameter $\beta_{N,p}$, a lattice configuration with a group element on
each oriented link, and an insertion $I_{N,p}(v,x;U)$ for each microscopic
field token $v$. The torus has $D_N L_p$ sites per side, and the physical mesh
is $a_N=D_N^{-1}$. The period and mesh schedules, action parameter and
insertions all belong to one family chosen before the requested IR sequence.
\src{O1,N3}

Writing $S_{N,p}$ for the resulting lattice correlations, the common limit
pattern is schematically
\begin{equation}\label{eq:limits}
 \begin{gathered}
 \exists C\quad\forall p:\N\to\N\text{ strictly increasing}\quad
 \exists q:\N\to\N\text{ strictly increasing}\quad\exists S,\\
 S_{p(q(k))}=\lim_{N\to\infty}S_{N,p(q(k))},\qquad
 S=\lim_{k\to\infty}S_{p(q(k))}.
 \end{gathered}
\end{equation}
The correlations and their test domains must be supplied to make this
schematic formula precise. The UV limit is a full limit along the chosen
cutoff schedule at each fixed period; the IR limit uses a subsequence of the
requested period indices. Neither $p$ nor $q(k)$ is itself the physical side
length. One refinement serves the whole hierarchy covered by the theorem,
rather than being reselected for each correlation function.

November's \code{FixedPointReturns} makes the intermediate objects explicit.
Before requesting an IR sequence, one fixes a table $S_p(v;x)$ of continuous
correlation kernels on separated configurations inside each physical torus.
There are two compatibility statements. Smeared lattice correlations tend
to integrals against that table. Also, on a compact separated set, lattice
moments converge uniformly when their insertion sites approximate the
specified continuum sites. The latter is an $\varepsilon$--$\delta$ statement
about the mesh and point positions, not just pointwise convergence at one
preferred sequence of lattice sites. \src{N4}

The \code{pointwise} field of November's \code{ContinuumLimit} then requires
these same finite-period kernels to converge locally uniformly along the
selected IR subsequence. Its \code{smeared} field identifies the same mean
hierarchy with the ordered lattice test limits. This exposes both steps of
\eqref{eq:limits} in the public certificate. It does not make the IR limit
unique and does not supply a general computable convergence rate. The
selection uses ordinary noncomputable existence machinery.

This distinction also prevents confusion between two uses of ``UV limit.''
Removing the lattice cutoff at fixed physical volume is one limit. Taking
short-distance scaling limits \emph{inside the already constructed continuum
theory} is another. October has full convergence for the first operation
even though the Maxwell comparison retained in our target uses a shrinking
subsequence for the second. Section~\ref{sec:maxwell} explains November's
improvement to the latter.

\section{How much of the common proof is actually unchanged?}
\label{sec:common}

There are two useful comparison records, at different levels. First, a
source-import analysis compares October's \code{Construction} with November's
\code{Comparators.Existence.Solution}. These closures contain 9,466 and
9,997 local modules respectively, excluding external dependencies. The
November closure has the following classification against the original
October source tree. \src{E2}

\begin{center}
\begin{tabular}{lr}
\toprule Classification & November existence modules\\\midrule
Byte-identical source & 8,130\\
Differences only in comments or whitespace & 959\\
Differences only in diagnostics after comment removal & 281\\
Remaining code differences at shared paths & 77\\
New paths & 550\\\midrule
Total & 9,997\\\bottomrule
\end{tabular}
\end{center}

These are whole imported modules, not a count of declarations required by
the final proof. Some new paths are reorganized or renamed old material.
Consequently the table measures reuse and helps direct review; it does not
say that precisely 550 new mathematical results were needed. Nor should the
size of the entire November repository, much of which concerns matching,
be used to estimate how much the common existence argument changed.

Second, a retained Lean diagnostic starts at \code{YM.Spec.Final} in each
compiled environment and follows project constants occurring in types and
definition bodies, including constructors of encountered inductives. It
does not follow theorem proof bodies. After removing binder names and
expression metadata, all 1,591 recorded declarations have identical
digests: declaration kinds, universe parameters, types and recorded bodies
agree, and neither side has an extra declaration in this closure. The
diagnostic preserves constant names rather than unfolding them into an
arbitrary equivalence class. The two digest inventories have been rechecked
for equality in preparing this report. \src{E2}

This is substantial evidence that the retained native certificate has not
quietly acquired a different meaning. It is not a theorem identifying
October's selected Wilson family with November's selected family, a replay
of every theorem proof, or an audit of all the new public November records.
The chosen witnesses may differ even when they satisfy the same native
predicate. New public obligations still require inspection on their own.

\section{What the November existence specification adds}
\label{sec:additions}

\subsection{Microscopic fields and the actual Wilson parameter}

November's \code{WilsonData} supplies finite-dimensional raw coefficient
spaces at each engineering cap. Their tokens exhaust the microscopic token
type, and their finite-lattice insertions are complex-linear in the raw
coefficients. There are compatible conjugation and reflected-adjoint maps,
translation covariance, and a unit insertion that equals one at sufficiently
fine cutoff, uniformly over periods. \code{AdjointOn} requires a surjective
projection onto physical fields, a linear quotient at each grade, and
compatibility with the reconstructed adjoint and unit. \src{N5}

These conditions address a concrete weakness of a target that allows any
local gauge-invariant observable. Setting every insertion to the constant
one obeys locality and gauge invariance, but violates linearity at zero.
Linearity alone would still allow uninformative field systems; coverage of
the physical fields, the graded Maxwell identification, nonzero trace law
and other requirements must be considered together. Our October target
instead retains native local-polynomial sections and their prescribed
classical curvature-jet expansion. These are related approaches, not an
already proved equivalence of specifications.

The one-loop coefficient is now expressed using the same trace-normalized
color space as the free reference and the Wilson action. If $E_a$ is an
orthonormal Lie-algebra frame for
$\langle X,Y\rangle=-m^{-1}\operatorname{ReTr}(XY)$, with $m$ the matrix
size and $d_{\g}=\dim_\R\g$, the public definition is
\begin{equation}\label{eq:casimir}
 C_{\g}=\frac{1}{m d_{\g}}\sum_{a,b}
       -\operatorname{ReTr}\bigl([E_a,E_b]^2\bigr),\qquad
 b_0=\frac{11C_{\g}}{48\pi^2}.
\end{equation}
The running clause requires $\beta_{N,p}/\log D_N\to 2b_0$ uniformly in
$p$. This is a statement about the parameter in the actual Wilson law.
\code{WilsonTrajectory} additionally fixes one reference inverse coupling
for all volumes, constrains all finite bare parameters to intervals around
its logarithmic center, and imposes divisibility, odd blocking factors and
a gentle relative logarithmic advance of the cutoff. These clauses prevent
the running law from referring to an independently chosen number unrelated
to the lattice action. They do not prove uniqueness of the IR state or the
existence of a particular intrinsic physical mass scale. \src{N5,N6}

There is also a useful uniqueness statement of limited scope:
\code{AdjointOn.orthogonalSection} chooses the unique lift perpendicular to
the kernel of the raw-to-physical map. This fixes a representative once the
raw coordinates and projection are given. It is not uniqueness of the
renormalized theory, its physical field normalization, or its OPE. \src{N7}

\subsection{Phases are tied to the mean Wilson hierarchy}

The small October public target relates the mean and proposed phase
correlations through a mixture identity. A first-moment identity by itself
does not establish that the proposed pieces are physical states. Formally,
one could average $A$ and $2M-A$ to obtain $M$ even when the second term has
no positivity. This is an illustration of the missing relation, not a
complete countermodel satisfying every October condition.

November requires more. For a separated word $w$ of microscopic observables,
let $S_\chi(w)$ be its reading in character $\chi$. The phase signatures must
separate characters. If $w_1,\ldots,w_k$ are translated apart by increasing
Euclidean-time gaps, their single concatenated reading in the mean hierarchy
must satisfy
\begin{equation}\label{eq:trains}
 \lim_{n\to\infty} S_{\rm mean}
   \bigl(w_1,\tau_n w_2,\ldots,\tau_{(k-1)n}w_k\bigr)
  =\int_X\prod_{j=1}^k S_\chi(w_j)\,d\nu(\chi).
\end{equation}
The formal statement also permits OS-adjoint words, giving the corresponding
complex-conjugate factors. The cutoff and volume limits have already been
taken before this separation limit. Thus the same Wilson hierarchy
determines all these joint phase moments, not just a barycenter. \src{N8}

Every character must also have a Wightman theory, a faithful field
representation, analytic continuation and joint tempered boundary
distributions. The selected theory and its tube function are identified
with the selected member of this family. These are genuine additional
requirements on the proposed phases. Our October target retains a native
spectral certificate that supplies related information, but the comparison
of the two formulations has not been packaged as a Lean implication.
Neither formulation classifies the vacuum space or proves that there is
only one character. \src{N9,G1}

\subsection{The anomaly field is the Ward-normalized stress trace}

The October compact hierarchy stores a field called \code{centeredTrace}
without storing a stress tensor in that structure. November stores
$T_{\mu\nu}$ and a unit field and defines
\[
 \Theta=\sum_\mu T_{\mu\mu}
       -S_{\chi_*}\!\left(\sum_\mu T_{\mu\mu}\right)\mathbf1.
\]
Its linearity and unit conditions make this subtraction agree with what
correlations and reconstructed operators actually read. \src{N10}

For a smooth compact window $\eta$ equal to one near the distinguished
insertion and supported away from the spectators, the translation Ward
identity requires
\begin{equation}\label{eq:ward}
 -\sum_\mu\int (\partial_\mu\eta)(y)
 S_s\bigl(T_{\mu\nu}(y),A(x_0),\ldots\bigr)\,d^4y
   =\partial_{x_0^\nu}S_s\bigl(A(x_0),\ldots\bigr).
\end{equation}
Here $s$ ranges over the mean and every phase. The same tensor has symmetry,
conservation, rotation covariance and rotation Ward laws. Its continued
Minkowski form is required to implement the infinitesimal symmetry action
on the reconstructed field domain. Equation~\eqref{eq:ward} fixes a sign
and normalization that conservation away from insertions alone would not
fix. A freely rescaled scalar cannot replace the trace while preserving
these relations. Our October Ward and trace bridges already check this
connection on the October native witness. \src{N11,N9,G1}

The explicit leading trace asymptotic remains
\[
 r^8\log^2(1/r)\,S_{\chi_*}\bigl(\Theta(-re_0/2),\Theta(re_0/2)\bigr)
 \longrightarrow \frac{3\dim_\R\g}{\pi^4}.
\]
The specified finite Wick-profile exclusion is retained. This should not
be enlarged to a theorem excluding every imaginable realization by
composites of an arbitrary Gaussian family. Also, the numerical $b_0$
cancels from this leading trace amplitude; the explicit running clause is
a distinct test. The group's worked October LeadingTrace example explains
that cancellation and its supporting normalization conventions. \src{G2,N10}

\subsection{An all-radius, all-phase Maxwell comparison}
\label{sec:maxwell}

The November free reference is tied to the chosen group. Its color space
maps onto the specified Lie algebra with the normalized trace metric.
Its variables are free curvature jets obeying the Bianchi and free field
equations; the sectors are exactly the gauge-invariant polynomials through
each engineering degree. Correlations are fixed by Wick pairing of the
massless curvature covariance, and the free stress has a specified
normalization. The reference is therefore constrained data rather than an
arbitrary correlation array to which an interacting theory may be fitted.
\src{N12}

A radius-dependent linear isomorphism $J_r$ identifies each such free
sector with the corresponding physical field space. For finitely many
fields and any derivative order, the public target requires, schematically,
\begin{equation}\label{eq:maxwell}
 \sup_{s\in X\cup\{\mathrm{mean}\}}\sup_{x\in K}
 \left\|\partial^\alpha_x
 S_s\bigl(J_rv_1,\ldots,J_rv_n;rx\bigr)
 -\partial^\alpha_x S_{\rm Maxwell}(v_1,\ldots,v_n;x)\right\|
 \longrightarrow0\quad(r\downarrow0)
\end{equation}
for compact separated $K$. The formal \code{mixed} clause also allows
marked slots containing the \emph{fixed physical stress tensor} multiplied
by $r^4$; these slots cannot be replaced by arbitrary moving fields.
Its eventual-radius quantifier precedes the phase quantifier, so the
comparison is uniform over phases. \src{N13}

The target also bounds $J_r$ and its inverse in fixed coordinates, with
arbitrarily small power losses. A row that first detects engineering degree
$e$ has a forward bound of order $r^{e-\delta}$, and the inverse on the
space through degree $d$ has a bound of order $r^{-d-\delta}$. Fixed-field
laws identify the unit, stress and derivatives and include a finite set of
separating probes. Such conditions restrict the otherwise substantial
freedom of a radius-dependent choice of basis.

Our October target explicitly retains \code{Pert.FirstExtension}, which
compares the mean hierarchy and its derivatives with the free hierarchy
along a common shrinking sequence, including marked stress slots. Its
statement alone does not contain November's all-radius extension, uniform
every-phase conclusion or frame bounds. November's proof obtains sampled
estimates and applies \code{AllRadii.extendFrame} and
\code{AllRadii.maxwellLimit} to supply the stronger public conclusion.
This is a substantive difference between the compared statements. A more
extensive theorem elsewhere in October could be relevant, but no completed
implication from our published combined target to this November clause is
claimed here. \src{O2,N1,N13}

Equation~\eqref{eq:maxwell} is an order-zero free limit with controlled
fields. It is not the deferred theorem identifying positive-order
perturbative coefficients.

\subsection{OPE consistency exposed in the existence target}

The common OPE subject remains an asymptotic expansion of separated local
fields with coefficient bounds and errors after finite truncation.
November adds \code{CoefficientLaws}: multilinearity, compatibility with
the physical labels and translations, dual channel coordinates, completeness
of the channel expansion of one field, and deletion of a unit insertion.
One explicit consequence is that a channel field cannot be zero; this
excludes padding the coefficient table with invisible zero fields. \src{N14}

\code{PowerTreeExpansion} also requires estimates for a fixed finite fusion
tree with internal truncation levels fixed before the outer radius tends
to zero. It compares flat and nested coefficients, and controls the
corresponding correlation remainder with spectators. The bounds permit
arbitrarily small power losses and include tested statements. These clauses
expose consistency and quantitative information beyond a bare existential
OPE assertion. They do not identify the coefficients with Feynman-graph
coefficients, prove a convergent infinite OPE series, or establish uniqueness
of every possible contact extension. The group's explicit-radius OPE
predicate is now proved on the November hierarchy (Section~\ref{sec:remaining}).
Equivalence between the finite-tree laws and October's native nested OPE
formulation remains a separate alignment question. \src{N15,G3}

\section{How the strengthened requirements are proved}
\label{sec:proofs}

The differences above are differences between specifications. To assess
what changed mathematically, one must also read the proof supplying each
new clause. Some clauses are direct translations of October's native
certificate; others need additional estimates and selection arguments.
The following comparison traces the actual November existence proof,
without using the separate perturbative matching theorem.

\subsection{One construction supplies all the objects}

November's \path{Comparators/Existence/Solution.lean} starts by constructing
one native run. Before any volume sequence is requested, it fixes the
Wilson family, its lattice fields, a free reference for the same gauge
group, and the finite-period continuum point table. For a requested volume
sequence it invokes the native refinement theorem, obtaining a table $S$,
a certificate $F$ of type \code{YM.Spec.Final}, and an OPE reader $Q$.
The public hierarchy is then literally \code{hierarchy F Q}. The proof
fills the fields of \code{ContinuumLimit} using these arguments throughout.
\src{N1}

This shared dependence is the essential connection. The phase proof does
not start from a second independently selected ensemble, nor does the Ward
proof select a new stress tensor after testing the anomaly. The OPE
coefficients are the coordinates of the reader $Q$ belonging to this
hierarchy. Its Maxwell comparison also uses $F$ and $Q$. This is visible in
the short final proof and in the types of its bridge lemmas; it is more
informative than a list of separate existence statements.

\begingroup\small\setlength{\tabcolsep}{4pt}
\renewcommand{\arraystretch}{1.15}
\begin{longtable}{@{}>{\raggedright\arraybackslash}p{.23\textwidth}>{\raggedright\arraybackslash}p{.40\textwidth}>{\raggedright\arraybackslash}p{.30\textwidth}@{}}
\toprule Public obligation & November producer & What is identified\\\midrule\endfirsthead
\toprule Public obligation & November producer & What is identified\\\midrule\endhead
Finite-period and IR points & \code{LeanPointReturns.nativeReturns} & One finite-period table works with every native final hierarchy of this run.\\
Microscopic adjoint and quotient & \code{IssuedBridge.FieldQuotient.adjoints} & The actual lattice token projects to the same physical field and reflected adjoint.\\
Phase moment laws & \code{IssuedBridge.PhaseCorrelations.law} & Word readings are characters of the native ground-state algebra.\\
Every-phase reconstruction & \code{VacuumReconstruction.phases} in \code{IssuedExistencePhase} & The existing phase atlas supplies the Wightman theories and boundary distributions.\\
Stress and anomaly & \code{publicWards}, \code{publicTrace} & The tensor is $F$'s energy field, and the anomaly field is its centered trace.\\
OPE algebra and trees & \code{publicCoefficientLaws}, \code{onFiledUniform} & Coefficients, nested contractions and spectator estimates use the same $Q$.\\
Full short-distance comparison & \code{FixedFields.sampledEstimates}, \code{AllRadii.extendFrame} & Sampled and all-radius field frames are reconciled on the same $F$.\\
\bottomrule
\end{longtable}\endgroup

These source-level identifications are now supplemented by the group's
combined November theorem, which requires the independent October checks
on the same native witness. This establishes that the November construction
satisfies those requirements; it does not prove equivalence between the
two acceptance specifications. Section~\ref{sec:remaining} describes the port.

\subsection{Point convergence is made explicit, using existing local control}

The proof of \code{nativeReturns} has two parts. At a fixed period, the
existing full-mesh UV theorem provides a continuous kernel and convergence
along every sequence of lattice sites approaching a separated continuum
configuration. Compactness of a fixed configuration set upgrades this to
the uniform $\varepsilon$--$\delta$ assertion in \code{FixedPointReturns}:
failure would give increasingly fine lattices and bad configurations;
a convergent subsequence of those configurations would contradict the
original point-limit assertion. This compactness argument does not replace
the UV limit by a subsequential one. It proves a uniform statement about
the already established full UV limit. \src{N16}

For the IR comparison, \code{SpatialModulus.purchased} obtains a local
Lipschitz estimate from the submitted bounds on lattice finite differences,
uniform in the relevant large periods and fine cutoffs. On a slightly
enlarged separated compact set, convergence at the existing dense set of
generic configurations, that Lipschitz estimate, and continuity of the
limiting kernel give local control near each point. A finite cover makes
the control uniform. Passing to the fixed-period UV limit then gives
local uniform convergence of the finite-period continuum kernels to $S$
along the chosen volume subsequence. \src{N16}

The November module packages this argument as a public bridge. The native
modules \path{YM/FullVolume/Profiles.lean} and
\path{YM/FullVolume/Connect.lean} used here are byte-identical to October's
copies. The improvement is an explicit point-table interface and the
uniform comparison proof connecting it to the final hierarchy, rather than
a new solution of the IR uniqueness problem.

\subsection{The phase and Ward checks expose native October structure}

For phase moments, November's \code{OriginalTrainMoments.original} uses
\code{YM.OS.Ground.arcs}. That lemma was already present in October, in
the byte-identical file \path{YM/Ground/Moments.lean}. Increasing the
separations between word insertions makes their mean correlation converge
to a product of compressed operators on the ground-state space. The
spectral representation turns this product into the integral of the
product of character readings in \eqref{eq:trains}. Adding the reflected
adjoint words gives the conjugate readings. \src{O3,N17}

The separation statement has a similarly concrete proof. The compressed
word operators generate the ground-state algebra. Two characters agreeing
on all word readings agree on these generators, hence on the algebraic
star algebra they generate, and then by continuity on its norm closure.
Injectivity of the spectral parametrization identifies the characters.
November exposes this consequence as injectivity of the public phase
signature. It supplies more than a choice of functions whose average
happens to be the mean correlation. \src{N17}

Every-phase reconstruction is also built from the native phase atlas
\code{L : EveryPhase.Almanac F}. Its existing actors provide the Hilbert
spaces and fields; its contours provide the analytic functions and joint
tempered boundary distributions. November checks the unit operator and
adjoint identifications and explicitly identifies the selected actor and
selected analytic function with the public selected theory. October's
detailed \code{constructed} theorem already retained this all-phase
reconstruction data. Thus much of the improvement here concerns what the
public Comparator target is required to retain. \src{O2,N18}

The Ward bridge is particularly transparent. For the mean hierarchy it
uses \code{F.phase.transport.charges}; for a character $\chi$ it uses
\code{F.phase.phase.flux chi}. Covariance likewise comes from
\code{transport.motion} and \code{phase.sym}. These are fields of the
native certificate whose definition closure agrees between the releases.
November's public checks now require the derivative insertion, tensor
transformation, translation flux, rotation flux and centered trace to fit
together. They expose structure already central to the October construction
and to our October Ward check. \src{N19,O4}

\subsection{All radii require a new selection argument}

One cannot obtain convergence at every small radius merely by writing down
values between the sampled radii of \code{FirstExtension}. November instead
proves that satisfactory frames exist at every sufficiently small radius
for each fixed finite collection of demands. A demand bounds the number of
insertions, filtration degrees and derivative orders, specifies a compact
separated configuration window, and fixes an error tolerance. The finite
tests are made on basis vectors; multilinearity then recovers general
fields. \src{N20}

The key lemma is \code{PhysMS.Frame.on_crew}. If a fixed demand failed at
arbitrarily small radii, one could select a bad sequence tending to zero.
The native trace-construction certificate supplies a refinement and
calibrated frames along that sequence. The sequence comparison eventually
satisfies the finite demand, giving a contradiction. The lemma therefore
establishes eventual existence at all real radii, using more information
than the single sampled sequence alone.

Next, \code{furnished} makes one radius-dependent choice satisfying an
increasing number of these demands as $r\downarrow0$.
\code{base_stage} assembles the resulting frame tower, with compatible
filtration inclusions, unit, reality and signed-coordinate symmetries, and
mean Maxwell convergence on every separated compact set. This is a
diagonal selection over finite requirements at each radius. It does not
assert that the chosen frame depends continuously or differentiably on
$r$, or give an effective convergence modulus. \src{N20}

Finally, the public bridge must also agree exactly with the original
sampled frames. \code{SparseFrame.splice} uses the original frame at a
sampled radius and the all-radius frame elsewhere. Both choices satisfy
the required asymptotics on their respective sets, so the combined choice
does too. The splice is a compatibility step after the all-radius theorem
has been proved; it is not the argument establishing convergence between
samples. \src{N21}

The native \code{PhysMS/Frame} modules and these assembly arguments are
new files in November. This is additional analysis supporting a stronger
existence endpoint. The source comparison does not prove that no equivalent
result could be derived from October's more extensive internal estimates.

\subsection{Why phases share the short-distance limit}

The phase comparison uses the OPE and the engineering filtration. For the
simple groups in the theorem, the free gauge-invariant field space below
degree four is the identity line. Equality of the finite-dimensional
physical and free field ranks transfers this statement to the interacting
field quotient. In \code{YM.OPE.Short.Gap}, \code{unit_field} proves that
every field of degree below four is a multiple of the unit, and
\code{no_channels} eliminates positive OPE channels of degrees one, two
and three. This is an engineering-degree statement, not a spectral mass
gap. \src{N22}

Consider a word $\boldsymbol A=(A_1,\ldots,A_n)$, with fixed filtration
degrees $d_i$, and rescale every insertion position by $r$. Truncating its
OPE at degree three leaves only the identity coefficient. The reader $Q$
is common to the phases, and each phase assigns expectation one to the
identity. Consequently the phase dependence is contained in the OPE
remainder. With $D=\sum_i d_i$, the internal estimate is
\begin{equation}\label{eq:phase-small}
 \sup_{s,t}\sup_{x\in K}
 \left\|\partial_x^j\left[
 r^D\bigl(S_s(\boldsymbol A;rx)-S_t(\boldsymbol A;rx)\bigr)
 \right]\right\|
 \leq 2(4n)^j r^{4-\varepsilon}
\end{equation}
at sufficiently small positive $r$, for a separated compact $K$ and
$\varepsilon>0$. Here $\partial_x^j$ denotes the full iterated derivative
in the rescaled coordinates, and $s,t$ may be characters or the mean.
For $j=0$ the numerical prefactor is two. Differentiating raises the
degree of the differentiated field and introduces the compensating power
of $r$, which preserves the remainder power. \src{N22}

The public \code{FixedFields.PhaseBound} deliberately asks for the weaker
zero-derivative bound $2r^{1-\delta}$, together with convergence of the
phase differences at every derivative order. The stronger
\eqref{eq:phase-small} is an internal theorem used by the existence proof;
it should not be attributed verbatim to the public target. Transferring
the estimate to moving fields also needs the frame bounds. The new
\code{PhysMS.Frame.Phases} argument supplies that transfer, obtaining
uniform vanishing of the differences in a common leading frame and then
combining it with the mean Maxwell limit. \src{N22,N23}

This explains both the strength and the limitation of the conclusion.
After normalization, the same free limit is reached in every phase, with
one frame shared by all phases. The frame covers the full finite-degree
gauge-invariant curvature-polynomial spaces, so this test is substantially
broader than testing only the stress trace. Yet all insertion points in
\eqref{eq:phase-small} are in the shrinking cluster. The argument does not
identify correlations at fixed macroscopic distances, nor assert the same
phase independence with additional unscaled spectator fields.

There is a useful consequence already available without reviewing
perturbative matching. For a calibrated coupling
$g(r)\asymp[\log(1/r)]^{-1/2}$, any positive power of $r$ is smaller than
any prescribed power $g(r)^N$ at sufficiently small $r$. The lemmas
\code{phase_jets_flat} and \code{Frame.negligible} make this comparison
precise. They show why different phases can have indistinguishable
short-distance power-series behavior while differing nonperturbatively.
They do not calculate that series or prove that its coefficients agree
with perturbation theory. \src{N22,N23}

\subsection{The fixed stress tensor survives the field-frame choice}

An arbitrary moving basis could otherwise conceal whether the physical
stress tensor has the required free limit. November treats it separately.
In \code{AllRadiiMaxwell.stress_column}, the rescaled physical tensor has
coordinates of the schematic form
\[
 U_r^{-1}(r^4T_{\mu\nu})
 =a(r)T^{\rm free}_{\mu\nu}
   +\delta_{\mu\nu}\bigl(d(r)\Phi+e(r)\mathbf1\bigr).
\]
The trace-calibration estimates give $a(r)\to1$ and $d(r)\to0$.
The unit term is controlled by the exact expectation identity
$4e(r)=r^4\langle\operatorname{tr}T\rangle$, so $e(r)\to0$ as well.
The free tensor here has the previously fixed Ward normalization.
\src{N24}

Finite-dimensional multilinearity then combines moving free-field slots
with these fixed stress slots. This supplies the marked mixed moments
in \code{MaxwellLimit.mixed}. The physical tensor is the one already
identified by the translation and rotation laws; it is not replaced by
a different tensor for each test or phase. The finite-probe conditions
in \code{FixedFieldLaws} also ensure that prescribed field columns can
be identified from a finite set of separated free-field measurements at
each degree. \src{N13,N24}

\subsection{Stronger existence checks do not re-export every October result}

November's public existence theorem is not literally a superset of the
long October \code{constructed} theorem. October explicitly returned
contact-extension devices and their compatibility with the OPE reader,
the \code{FourMeters} result, and a construction for every positive
construction-penalty parameter. November's short existence target has no
corresponding contact or \code{FourMeters} clause; its displayed proof
chooses the penalty to be one. It asks for refinements of every requested
volume sequence with full fixed-volume UV convergence. \src{O2,N1,N2}

These are differences in what the endpoints retain, not evidence that
November's source has lost the older capabilities. The group's November
port now recovers these additional conclusions on its native witness,
including the contact devices, \code{FourMeters}, and arbitrary positive
penalty. They are proved by explicit adapters using November's lemmas;
they do not follow merely from a pass of its short public target. \src{G3}

\section{November passes the group's October checks}
\label{sec:remaining}

\textbf{The November construction satisfies the group's October acceptance
conditions.} We proved this in Lean using the November construction results,
and both combined targets passed Comparator on 11 November. This includes
the independent OS comparison, the lattice and continuum-limit checks,
the translation Ward identity, operator content, the leading stress-trace
two-point law, the existing scale checks, and the OPE check. \src{G3,E4}

To apply the October checks, we adapted their interfaces to the November code.
The central lattice, OS, translation Ward, operator-content and leading-trace
predicates retain their statement bodies. The OPE check retains its radius
and state quantifiers, with the renamed public channel field. The additional
source conclusions bundled into the October target were recovered through
explicit adapters. The independent acceptance requirements were retained.
\src{G3,E4}

Two theorems make the result precise. The theorem
\code{YMReview.Acceptance.combined_acceptance} proves the ported October
combined acceptance target, retaining its quantifiers over input extraction
orders and positive construction penalties. The theorem
\code{YMReview.November.combined_acceptance} additionally requires November's
full public existence conditions and the group's checks to hold together.
It retains one Wilson family and, for each requested volume sequence, a
refinement with native data $F$ and $Q$; the public hierarchy is explicitly
\code{YangMills.hierarchy F Q}. Thus the tests apply to the same constructed
theory. Both proofs are in \path{checks/NovemberCombined.lean}. \src{G3,E4}

For the OS comparison, the port identifies the public centered stress trace
with the native field used in our bridge, then proves equality of their
Schwinger functions. Thus E0--E4 and E0$'$ apply to the selected scalar
hierarchy on the same diagonal-flat joint Schwartz tests as in October.
This does not establish unrestricted Euclidean contact-term positivity or
uniqueness of a contact extension. The additional source contact and
calibration conclusions are retained separately; bundling them does not
independently validate every source definition.

\begingroup\small\setlength{\tabcolsep}{4pt}
\renewcommand{\arraystretch}{1.15}
\begin{longtable}{@{}>{\raggedright\arraybackslash}p{.22\textwidth}>{\raggedright\arraybackslash}p{.43\textwidth}>{\raggedright\arraybackslash}p{.28\textwidth}@{}}
\toprule Check & Result on the November construction & Scope or remaining work\\\midrule\endfirsthead
\toprule Check & Result on the November construction & Scope or remaining work\\\midrule\endhead
Common native core & Both targets use a retained native \code{Final}; the new target also contains all November public existence conditions. & No cross-release equality theorem.\\
Independent OS axioms & E0--E4 and E0$'$ proved for the public selected centered stress trace. & Scalar hierarchy; diagonal-flat joint tests.\\
Lattice law and fields & Independent Wilson-law and insertion requirements proved on the November witness. & Predicate bodies retained.\\
Stress and leading trace & Independent Ward identity and leading two-point law proved, with the public trace identified. & Same sign, coefficient and field requirements.\\
Operator content & Independent coverage, quotient, no-extra-null-field and dense-domain requirements proved. & No separately specified canonical Hilbert space.\\
Common physical scale & Existing common calibration and positive-probe conclusions transferred. & No intrinsic nonzero mass scale or mass gap claimed.\\
OPE and contacts & Explicit-radius OPE check and additional contact/calibration source conclusions proved. & No contact-extension uniqueness or positive-order matching claim.\\
\bottomrule
\end{longtable}\endgroup

The scope is the same as for the October checks: an existence theorem with
suitable refinements of the requested volume sequences. The result does
not require October and November to select identical IR limits. A mass gap,
unique IR limit, unique vacuum and universality under arbitrary regulator
choices were not October acceptance conditions and are not established by
this transfer. Human validation of the definitions, a new November-specific
adversarial exercise and the review of positive-order perturbative matching
remain separate tasks; they do not change the completed status of the
formal October checks.

\section{Completed Comparator runs and their meaning}
\label{sec:verification}

The received November sources were rebuilt with Lean 4.32.0, beginning with
no compiled project objects reused from either submission. Verified external
dependency caches were reused. All four default targets compiled. The first
resource-timing wrapper recorded an operating-system reporting failure after
Lake had completed successfully; a separate direct
\code{lake --no-cache --no-build build} then returned zero. That distinction
is preserved in the receipts rather than replacing the original failed
wrapper record. \src{E1}

Each target was checked against its supplied challenge. Both configurations
permit only \code{propext}, \code{Classical.choice} and \code{Quot.sound},
enable no definition holes, and disable the optional second kernel. The
separate theorem-specific axiom audits report exactly those three axioms.
The challenge files intentionally use \code{sorry} to state the reference
obligation; an accepted solution cannot depend on \code{sorryAx}.

\begin{center}\small
\begin{tabular}{llll}
\toprule Target & Checker & Completion (UTC, 10 Nov) & Result\\\midrule
Existence & Release-pinned \code{07bc4ea} & 04:48:38 & Passed\\
Matching & Release-pinned \code{07bc4ea} & 05:30:13 & Passed\\
Existence & \code{d03acab}, Lean 4.34.0 & 05:43:19 & Passed\\
Matching & \code{d03acab}, Lean 4.34.0 & 08:18:57 & Passed\\\bottomrule
\end{tabular}
\end{center}

The group's two combined November targets passed together on 11 November
at 23:33:32 UTC, using Comparator \code{d03acab} and its Lean 4.34.0
checking kernel, with no definition holes. The run took 1,020 seconds and
built the review modules in a fresh overlay while reusing the existing
November and dependency builds. The separate compilation audit covered 49
review modules with no warnings and 95 axiom reports, all limited to the
three standard axioms. The original 37,673 submission files and the frozen
review inputs were unchanged across the run. This is an additional check
of the transferred October requirements, distinct from the supplied-target
runs above. \src{E4}

The repeat used an unchanged upstream Comparator checkout and a Lean 4.32.0
exporter, with the newer checking kernel. The local name ``patched recheck''
in its evidence filenames denotes this newer-tool replay; it does not mean
we edited the submitted proof or maintained a locally modified checker.
The newer replay matters because the local tool audit reproduced defects
with the older tools. This finding concerns the verifier software and is
distinct from the October target's omitted physical relations.

Preparing this report rechecked the four log hashes, exit statuses and the
explicit kernel/Comparator acceptance messages. It also rehashed all 37,673
files in the received November source manifest: no differences were found.
The prior four small controls accepted a valid proof and rejected a changed
statement, a forbidden axiom and a changed definition. Those controls are
not complete attempts to substitute a physically different QFT.

The runs used the upstream macOS development launcher
\code{fake-landrun.sh}. They establish the recorded formal comparison and
kernel verdicts under that execution mode. They do not reproduce Linux
process isolation, constitute a Palomar run, or validate the physical
meaning of perturbative matching. The continued scientific review should
use the stronger November statement as its object, not treat these mechanical
verdicts as a substitute for that review.

\appendix
\section{Source map and retained evidence}
\label{sec:sources}

October source links below point to the public original-submission commit
\code{3456e5e75cf8d4cd558955a5c7638fba2d68f464}; group-check links point to
the published YM assembly \code{55ba9fc36293b3f4d52c6b1f813169a0be76ee98}.
November source paths are relative to the \code{lean/} directory of the received
November 4 distribution. The group's port paths in G3 and E4 refer to its
separate November verification package. This report does not require a link
to a private GitHub repository. The SHA-256 of its supplied
\code{SOURCES.sha256} file is
\begin{quote}\small\path{c15ffe5e8ca808f5748170c2f3910a6000cf4a01c0757ece91bd5b5bf65e60f9}.\end{quote}

\begin{description}[style=nextline,leftmargin=3.5em,labelwidth=3em,align=left]
\item[O1] \oldlink{ChallengeSpec.lean}{\code{ChallengeSpec.lean}}: public regularity, hierarchy, trace test, properties and joint existence. \oldlink{PublicWilson.lean}{\code{PublicWilson.lean}} and \oldlink{PublicTheory.lean}{\code{PublicTheory.lean}}: lattice and Wightman definitions.
\item[O2] \oldlink{YM/Pert/FirstTarget.lean}{\code{YM/Pert/FirstTarget.lean}}, especially \code{FirstExtension}; \oldlink{Construction.lean}{\code{Construction.lean}}: the detailed October construction.
\item[O3] \oldlink{YM/Ground/Moments.lean}{\code{YM/Ground/Moments.lean}}, \code{YM.OS.Ground.arcs}: separated-word limits in the ground-state algebra; byte-identical in November.
\item[O4] \oldlink{YM/Spec/Closed.lean}{\code{YM/Spec/Closed.lean}}, \code{Locals}, \code{Materializes}, \code{Final}: the retained native field, Ward, phase and reconstruction data.
\item[G1] \reviewlink{checks/AcceptanceSpec.lean}{\code{checks/AcceptanceSpec.lean}}: \code{SourceConclusions}, \code{Accepted}, \code{CombinedAcceptance}; \reviewlink{checks/CombinedAcceptance.lean}{\code{checks/CombinedAcceptance.lean}} and \reviewlink{checks/CompleteOS.lean}{\code{checks/CompleteOS.lean}}: compiled group endpoints.
\item[G2] \reviewlink{docs/examples/LeadingTrace-20261028.lean}{\code{LeadingTrace-20261028.lean}} and \reviewlink{docs/examples/OctoberChain.lean}{\code{OctoberChain.lean}}: worked derivation and construction connections.
\item[G3] In the November verification package: \path{checks/NovemberAcceptanceSpec.lean} and \path{checks/NovemberCombined.lean}, the two combined targets and their proofs; \path{checks/NovemberOS.lean}, the public-trace and Schwinger-function identifications; \path{checks/RetainedSourceProof.lean} and \path{checks/RetainedContactProof.lean}, the additional source conclusions. \path{docs/november-port.md} gives the detailed port report and reproduction commands.
\item[N1] \path{Comparators/Existence/Challenge.lean}; \path{Comparators/Existence/Solution.lean}: the exact challenge and its construction proof.
\item[N2] \path{YangMills/Existence.lean}, \code{ContinuumLimit} and \code{Construction.Realizes}; \path{YangMills/VolumeLimit.lean}, \code{Construction}.
\item[N3] \path{YangMills/Geometry.lean}; \path{YangMills/Wightman.lean}; \path{YangMills/WilsonLattice.lean}.
\item[N4] \path{YangMills/FiniteVolume.lean}, \code{PeriodPointTable} and \code{FixedPointReturns}.
\item[N5] \path{YangMills/LatticeFields.lean}, \code{wilsonOneLoop}, \code{GradedTokens}, \code{WilsonData}, \code{AdjointOn}.
\item[N6] \path{YangMills/CutoffTrajectory.lean}, \code{WilsonTrajectory}.
\item[N7] \path{YangMills/LatticeSections.lean}, \code{AdjointOn.orthogonalSection}.
\item[N8] \path{YangMills/Phases.lean}, \code{ObservablePhases.Law}.
\item[N9] \path{YangMills/VacuumReconstruction.lean}, \code{TableLaws}, \code{ContinuedStress}, \code{PhaseSource}, \code{Phases}.
\item[N10] \path{YangMills/CorrelationHierarchy.lean}, \code{stressTrace}, \code{centeredTrace}, \code{TraceTest}.
\item[N11] \path{YangMills/StressTensor.lean}, \code{translationFlux}, \code{StressWardLaws}.
\item[N12] \path{YangMills/FreeFields.lean}, \code{Coordinates}, \code{Reference}.
\item[N13] \path{YangMills/ShortDistance.lean}, \code{FixedFieldLaws}, \code{FrameBounds}, \code{MaxwellLimit}; \path{Solution/AllRadiiBridge.lean}: public all-radius bridge.
\item[N14] \path{YangMills/OPECoefficientLaws.lean}, including \code{CoefficientLaws.channelField_ne_zero}.
\item[N15] \path{YangMills/PhysicalOPE/PowerTrees.lean}, \code{PowerTreeExpansion}.
\item[N16] \path{Solution/PointReturns.lean}, \code{SpatialModulus.purchased}, \code{CompactTransfer.measured}, \code{LeanPointReturns.Uniform.compact}, \code{LeanPointReturns.nativeReturns}: fixed-period point tables and local uniform IR comparison.
\item[N17] \path{YangMills/WilsonProof/NativeSupport/IssuedBridge/OriginalTrainMoments.lean}, \code{original}; \path{YangMills/WilsonProof/NativeSupport/IssuedBridge/OriginalPhaseSeparation.lean}, \code{original}: native train moments and character separation.
\item[N18] \path{Solution/PhaseReconstruction.lean}, \code{OriginalPhaseContours.source}, \code{VacuumReconstruction.phases} in \code{IssuedExistencePhase}: all-phase public reconstruction.
\item[N19] \path{YangMills/WilsonProof/NativeSupport/WardBridge.lean}, \code{publicWards}: public stress laws from the native mean and phase certificates.
\item[N20] \path{YM/PhysMS/Frame/Atlas.lean}, \code{survey}; \path{YM/PhysMS/Frame/Stage.lean}, \code{on_crew}, \code{furnished}, \code{base_stage}: finite demands and the all-radius frame selection.
\item[N21] \path{Solution/SelectedJets.lean}, \code{SparseFrame.splice}, \code{SparseFrame.original}, \code{SparseLeading.original} in \code{YM.Existence.SelectedJets}: compatibility with the sampled frames.
\item[N22] \path{YM/OPE/Short/Gap.lean}, \code{unit_field}, \code{no_channels}; \path{YM/PhysMS/Foundation/Readout.lean}, \code{phase_values}, \code{phase_jets}; \path{YM/PhysMS/Foundation/Estimates.lean}, \code{phase_jets_flat}: engineering-degree and phase-remainder estimates.
\item[N23] \path{YangMills/FieldMeasurements.lean}, \code{PhaseBound}; \path{YM/PhysMS/Frame/Phases.lean}, \code{within_price}, \code{negligible}, \code{all_leading}: public phase bound and its stronger native moving-frame estimates.
\item[N24] \path{Solution/AllRadiiMaxwell.lean}, \code{stress_column}, \code{allRadii}, \code{Law.marked}; \path{Solution/FixedFieldBridge.lean}, \code{detector}, \code{sampledEstimates}: fixed stress and finite-probe identification.
\end{description}

The retained local evidence uses the workspace root
\begin{quote}\small\path{/Users/oai/Documents/Codex/2026-10-28/thre/}.\end{quote}
The following paths are relative to that root:
\begin{description}[style=nextline,leftmargin=3.5em,labelwidth=3em,align=left]
\item[E1] \path{work/ym-nov4-verification/evidence/}: \code{comparator-results.json}, \code{patched-comparator-recheck.json}, the four corresponding Comparator logs and receipts, \code{release-build-check.json}, \code{prepared-inputs.json}, \code{controls.json}, and \code{verification-report.md}. The prepared source is \path{work/ym-nov4-verification/release/}.
\item[E2] \path{work/ym-october-november-comparison/}: \code{existence-only-import-comparison.json}, \code{public-declaration-comparison.json}, \code{DefinitionAudit.lean}, \code{definition-audit-receipt.json}, \code{definition-comparison.json} and the two definition-digest inventories. The older broad narrative report in that directory contains a stale pending-run sentence; this report uses the completed receipts in E1 for status.
\item[E3] \path{work/ym-october-november-acceptance/evidence/status.json}.
\par This records an earlier, stopped attempt to apply November tests to the October submission, with no Comparator start or successful complete combined build. It is separate from the successful application of the group's October checks to November recorded in E4.
\item[E4] In \path{/Users/oai/Documents/Codex/repos/YM-lean/OpenAI-Nov4-verification/}, at revision \code{7a15e748eedc89e951be4dbc0ce5200478171451}: \path{evidence/independent-checks/AxiomAudit.json} and \path{port-audit.json}; \path{evidence/combined-targets/combined-comparator.json}, its log, \path{prepared-inputs.json}, and \path{input-stability.json}. The frozen reference statements and configuration are in \path{comparator/NovemberChallenge.lean} and \path{comparator/combined.json}.
\end{description}

The adjacent \code{evidence-summary.json} records the source and log checks
performed for this report. \code{proof-path-audit.json} records the file
hashes and declaration locations used for Section~\ref{sec:proofs}, including
the byte-identical native modules identified there.
The new independent-target run is recorded separately in E4. A later review of perturbative
matching should start from its precise public reference, coefficient
definitions and same-witness clauses, rather than extending the conclusions
of this existence comparison by implication.
\end{document}
